% =============================== PORTADA ===============================
\thispagestyle{empty}
\setcounter{page}{0}
\vspace*{-1.2cm}
\noindent\begin{tikzpicture}[remember picture,overlay]
  \fill[navy] (current page.north west) rectangle ([yshift=-2.4mm]current page.north east);
  \fill[acento] ([yshift=-2.4mm]current page.north west) rectangle ([yshift=-3.0mm]current page.north east);
  \fill[navy] (current page.south west) rectangle ([yshift=2.4mm]current page.south east);
  \fill[acento] ([yshift=2.4mm]current page.south west) rectangle ([yshift=3.0mm]current page.south east);
\end{tikzpicture}%
\begin{tikzpicture}
  \useasboundingbox (0,-1.15) rectangle (\textwidth,0.75);
  \node[anchor=west,inner sep=0pt] at (0,0) {\includegraphics[width=0.22\textwidth]{logo-usal.png}};
  \node[anchor=east,inner sep=0pt] at (\textwidth,0) {\includegraphics[width=0.15\textwidth]{logo-ciencias.png}};
  \node[align=center,inner sep=0pt] at (0.535\textwidth,0)
    {{\color{navy}\large\espaciado[150]{Universidad de Salamanca}}\\[4pt]
     {\color{acento}\espaciado[110]{Facultad de Ciencias}\enspace\textperiodcentered\enspace\espaciado[110]{Grado en Física}}};
  \draw[regla,line width=0.6pt] (0,-1.1) -- (\textwidth,-1.1);
\end{tikzpicture}

\vspace{0.7cm}
\noindent{\color{acento}\large\espaciado[170]{Termodinámica I}\enspace\textperiodcentered\enspace\espaciado[170]{Tema 3}}\\[10pt]
\noindent{\color{navy}\fontsize{38}{44}\selectfont\bfseries Segunda ley}\\[8pt]
\noindent{\color{navy!80!black}\Large\itshape Entropía, reversibilidad y límites del trabajo}\\[10pt]
\noindent{\color{navy}\rule{\textwidth}{1pt}}\\[-8pt]
\noindent{\color{acento}\rule{\textwidth}{0.35pt}}

\vspace{0.35cm}
\noindent
\begin{minipage}[t]{0.32\textwidth}
  {\color{acento}\footnotesize\espaciado[180]{Autor}}\\[1pt]
  Daniel Martín Guerrero
\end{minipage}\hfill
\begin{minipage}[t]{0.32\textwidth}
  {\color{acento}\footnotesize\espaciado[180]{Profesor}}\\[1pt]
  Julián González Ayala
\end{minipage}\hfill
\begin{minipage}[t]{0.32\textwidth}\raggedleft
  {\color{acento}\footnotesize\espaciado[180]{Curso}}\\[1pt]
  2026/2027
\end{minipage}

\vspace{0.3cm}
\begin{center}
\begin{tikzpicture}
\begin{groupplot}[group style={group size=2 by 1,horizontal sep=1.5cm},art,width=7.3cm,height=5.0cm]
% ---- plano p-V con familias de isotermas y adiabáticas (gamma = 5/3)
\nextgroupplot[xmin=0.5,xmax=5.2,ymin=0,ymax=1.75,xtick=\empty,ytick=\empty,xlabel={$V$},ylabel={$p$}]
  \foreach \T in {0.7,1.0,1.5,2.1}{
    \edef\tmp{\noexpand\addplot[regla,line width=0.45pt,domain={max(0.5,\T/1.75)}:5.2,samples=60] {\T/x};}\tmp }
  \foreach \K in {0.8,1.5,2.763,4.3}{
    \edef\tmp{\noexpand\addplot[regla,line width=0.45pt,densely dashed,domain={max(0.5,(\K/1.75)^(0.6))}:5.2,samples=60] {\K*x^(-5/3)};}\tmp }
  \fill[fondomedio,opacity=0.9] plot[domain=1:2.5,samples=50] (\x,{1.5/\x})
    -- plot[domain=2.5:4.593,samples=60] (\x,{0.6*(2.5/\x)^(5/3)})
    -- plot[domain=4.593:1.837,samples=60] (\x,{1/\x})
    -- plot[domain=1.837:1,samples=60] (\x,{1.5*(1/\x)^(5/3)}) -- cycle;
  \addplot[rojo,medio=0.5,domain=0:1,samples=60] ({1+1.5*x},{1.5/(1+1.5*x)});
  \addplot[navy,medio=0.5,domain=0:1,samples=60] ({2.5+2.093*x},{0.6*(2.5/(2.5+2.093*x))^(5/3)});
  \addplot[acento,medio=0.5,domain=0:1,samples=60] ({4.593-2.756*x},{1/(4.593-2.756*x)});
  \addplot[navy,medio=0.55,domain=0:1,samples=60] ({1.837-0.837*x},{1.5*(1/(1.837-0.837*x))^(5/3)});
  \estado{1,1.5}{south east}{1}
  \estado{2.5,0.6}{south west}{2}
  \estado{4.593,0.21772}{south west}{3}
  \estado{1.837,0.54433}{north east}{4}
% ---- plano T-S
\nextgroupplot[xmin=0.4,xmax=3.0,ymin=0,ymax=1.95,xtick={1,2.4},xticklabels={$S_1$,$S_2$},
    ytick={1,1.5},yticklabels={$T_c$,$T_h$},xlabel={$S$},ylabel={$T$}]
  \foreach \y in {0.5,0.75,1.25,1.75}{\edef\tmp{\noexpand\draw[regla,line width=0.45pt] (axis cs:0.4,\y) -- (axis cs:3.0,\y);}\tmp}
  \foreach \x in {0.7,1.7,2.7}{\edef\tmp{\noexpand\draw[regla,line width=0.45pt,densely dashed] (axis cs:\x,0) -- (axis cs:\x,1.95);}\tmp}
  \fill[pattern=north east lines,pattern color=acentoclaro] (axis cs:1,0) rectangle (axis cs:2.4,1);
  \fill[fondomedio] (axis cs:1,1) rectangle (axis cs:2.4,1.5);
  \addplot[rojo,medio=0.5] coordinates {(1,1.5) (2.4,1.5)};
  \addplot[navy,medio=0.55] coordinates {(2.4,1.5) (2.4,1)};
  \addplot[acento,medio=0.5] coordinates {(2.4,1) (1,1)};
  \addplot[navy,medio=0.55] coordinates {(1,1) (1,1.5)};
  \node[font=\scriptsize,navy] at (axis cs:1.7,1.25) {$|W|$};
  \node[font=\scriptsize,acento,fill=white,inner sep=1.5pt] at (axis cs:1.7,0.5) {$|Q_c|$};
  \estado{1,1.5}{south east}{1}
  \estado{2.4,1.5}{south west}{2}
  \estado{2.4,1}{north west}{3}
  \estado{1,1}{north east}{4}
\end{groupplot}
\end{tikzpicture}\\[2pt]
{\footnotesize\itshape\color{gris}El ciclo de Carnot en los planos del trabajo y del calor, sobre las familias de isotermas
(continuas) y adiabáticas reversibles (discontinuas) del gas perfecto.}
\end{center}

\vspace{0.3cm}
\noindent{\color{regla}\rule{\textwidth}{0.6pt}}\\[2pt]
{\color{acento}\espaciado[200]{Resumen}}\par\smallskip
\noindent{\small El calor reversible admite el factor integrante $1/T$: $\dd S=\dbar Q_{\mathrm{rev}}/T$ define la
entropía, una función de estado extensiva. La desigualdad de Clausius y el balance
$\Delta S=\int\dbar Q/T+S_{\mathrm{gen}}$, con $S_{\mathrm{gen}}\geq0$, contienen la segunda ley, cuyos enunciados de
Clausius, Kelvin--Planck, Carnot y del aumento de la entropía son equivalentes. Con ellos se calculan rendimientos
de motores, refrigeradores y bombas de calor, se acota todo motor por $\eta_C=1-T_c/T_h$ y se obtiene el trabajo
máximo. Los seminarios resuelven, paso a paso y con todo el desarrollo simbólico, los nueve problemas del tema y las
cuestiones del parcial que corresponden a la segunda ley.\par}
\smallskip
\noindent{\footnotesize{\color{acento}\espaciado[200]{Palabras clave}}\enspace
máquina térmica\sepdot rendimiento y COP\sepdot enunciados de Clausius y Kelvin--Planck\sepdot entropía\sepdot
desigualdad de Clausius\sepdot producción de entropía\sepdot ciclo de Carnot\sepdot ecuaciones $T\,\dd S$\sepdot
ciclos Otto y Brayton\sepdot trabajo máximo\par
\smallskip
\noindent{\color{acento}\espaciado[200]{Clasificación PACS}}\enspace
05.70.$-$a (Termodinámica), 05.70.Ce (Funciones termodinámicas), 07.20.Pe (Máquinas térmicas)\par}

\vfill
\noindent{\color{regla}\rule{\textwidth}{0.6pt}}%
\vspace{-6pt}
\[
  \dd S=\frac{\dbar Q_{\mathrm{rev}}}{T}\qquad
  \Delta S_{\mathrm{univ}}=S_{\mathrm{gen}}\geq0\qquad
  \eta=1+\frac{Q_c}{Q_h}\leq1-\frac{T_c}{T_h}\qquad
  \dd U=T\,\dd S-p\,\dd V
\]\par
\vspace{-8pt}
\noindent{\color{regla}\rule{\textwidth}{0.6pt}}%

\vspace{0.25cm}
\begin{center}
  {\color{navy}\large\espaciado[220]{Apuntes formalizados con seminarios prácticos}}\\[4pt]
  {\color{acento}\espaciado[200]{Grado en Física\enspace\textperiodcentered\enspace Salamanca, curso 2026/2027}}
\end{center}
\vspace*{0.1cm}
